%_____________________________________________________________________________________________ 
% Chapter 7: Conclusion
%_____________________________________________________________________________________________ 
\chapter{Conclusion}
\label{sec:conclusion}

We have presented the design and evaluation of \persisthotstuff{}, a modular framework that
closes four practical gaps in existing HotStuff prototypes: crash-safe persistence, dynamic
membership, pluggable application logic, and liveness under low load.

All four extensions are fully implemented and rigorously tested in Rust.  Beyond simulation
and stress testing, we provide two new contributions that strengthen the technical
foundation:

\begin{itemize}[itemsep=2pt]
  \item A TLA+ specification of the dynamic membership protocol, verified by TLC across
    7.3\,M~distinct states with zero violations of 11~safety invariants; including quorum
    intersection, commit safety across epochs, and valid transition enforcement.
  \item A performance benchmark at cluster sizes $n \in \{4, 7, 10, 13\}$ confirming $O(n)$
    message complexity ($3(n{-}1)$ messages/round, matching the analytical HotStuff bound)
    and demonstrating $9.0\times$ fewer messages than PBFT at $n{=}13$.
\end{itemize}

The implementation is available on GitHub~\cite{rust-implementation}.

If there is one takeaway, it is that the gap between a BFT \emph{protocol} and a BFT
\emph{system} is larger than it looks.  Storage, membership, application interfaces, and
liveness under real workloads all matter, and addressing them together rather than one at a
time forces design decisions that would not surface otherwise.  We hope this blueprint is
useful to others working in the same space.
